\documentclass[11pt]{article}
\usepackage[a4paper,margin=1in]{geometry}
\usepackage[T1]{fontenc}
\usepackage{amsmath,amssymb,booktabs,tabularx}
\usepackage[authoryear,round]{natbib}
\usepackage[hidelinks]{hyperref}
\usepackage{siunitx}
\sisetup{group-digits=integer,group-separator={,},group-minimum-digits=4}
\newcommand{\conc}{\mathrm{C}}
\newcommand{\spr}{\mathrm{S}}
\title{Patient Shortage in Histopathology:\\Classifier or Coverage Limitation?}
\author{Yannik Queisler}
\date{September 10, 2026}

\begin{document}
\raggedbottom
\setlength{\emergencystretch}{2em}
\maketitle
\begin{abstract}
\noindent Patient shortage can impair histopathology classification even when class patch counts remain fixed. This experiment tests whether changing the classifier recovers the lost performance or whether broader training-patient coverage remains beneficial. The original multilayer perceptron (MLP), regularized multinomial logistic regression, and cosine nearest neighbours are compared on identical frozen Virchow2 features and matched concentrated and spread allocations across three patient-disjoint splits of BRACS and TCGA-UT. Under concentrated support, logistic regression improves patient-macro balanced accuracy over the MLP by only 0.19 percentage points on BRACS (95\% interval: $-0.59$ to 1.01) and 0.48 points on TCGA-UT (0.28 to 0.70). Broader coverage improves logistic regression by 2.44 and 4.70 points, respectively. Nearest neighbours underperform the MLP in both conditions and datasets. These results favour a coverage limitation within the tested classifiers, with a small additional classifier gain on TCGA-UT that falls below the prespecified one-point practical threshold. They do not establish that other classifiers could not recover more of the shortage deficit.
\end{abstract}

\tableofcontents
\clearpage

\section{Motivation}
\label{sec:question}
The preceding benchmark found that independent-support deprivation reduced balanced accuracy by \num{2.62}--\num{3.03} percentage points on BRACS and \num{4.76}--\num{4.83} points on TCGA-UT, despite matched class patch counts. No exploratory method produced a positive BRACS recovery interval excluding zero; the strongest TCGA-UT recoveries closed only about one sixth of the deficit \citep{queisler2026benchmarkResults}. These findings establish damage that the tested interventions largely failed to repair. They do not establish that repair is impossible.

The experiment distinguishes two explanations for this loss. Under a \emph{classifier limitation}, the frozen features still contain useful predictive structure, but the original multilayer perceptron (MLP) does not exploit it well with the available training data. A different classifier could therefore improve accuracy using exactly the same shortage allocation. Under a \emph{coverage limitation}, the training set contains too few distinct patients to represent the variation encountered in unseen patients. Broader patient coverage should then benefit several classifiers, even if changing the classifier alone offers little improvement. Both limitations can contribute.

The encoder maps each patch \(x\) to a fixed feature vector \(z=\operatorname{Virchow2}(x)\). Reducing the training-patient set changes the labelled examples available to learn a classifier; the features of a given test patch remain fixed. The central question is therefore whether better decisions can be learned from the shortage data, and how much additional patient coverage helps. This report calls this ability \emph{classification generalization to unseen patients}: how successfully a classifier trained on the available frozen features predicts classes in held-out patients.

Two alternative \emph{readouts}, or classifiers applied to the frozen features, complement the original MLP. A linear probe tests whether a linear decision rule can use the available class structure \citep{alain2016probes}. A nearest-neighbour classifier tests whether nearby feature vectors have informative labels; this also provides an established way to evaluate frozen representations \citep{caron2021emerging}. Their gains under shortage and their responses to broader coverage distinguish the two explanations within the tested classifier families.
\section{Controlled comparison}
\label{sec:design}
The experiment retains BRACS's seven lesion classes and TCGA-UT's 30 eligible cancer types, including the benchmark's exclusions and label order. Each patch retains its frozen 2560-dimensional Virchow2 feature vector, formed by concatenating the CLS token and mean patch token. All eligibility rules and the three existing train--validation--test splits remain fixed \citep{queisler2026benchmarkProtocol}. A partitioning unit is a BRACS patient or TCGA-UT participant; both are called patients below. Patients are disjoint between training, validation, and test within each split. Repeated splits can contain the same patient in different roles.

The comparison uses balanced class allocations. Concentrated support (\(\conc\)) draws the available patches from fewer patients; spread support (\(\spr\)) uses broader training-patient coverage with the same class-specific patch counts. Both conditions retain the benchmark's exact patch identities: the shared concentrated-balanced allocation and the spread allocation under the native class assignment.

\begin{table}[htbp]
\centering
\small
\renewcommand{\arraystretch}{1.15}
\begin{tabularx}{\textwidth}{@{}p{0.22\textwidth}XX@{}}
\toprule
Design element & Fixed or varied & Purpose \\
\midrule
Datasets and splits & BRACS and TCGA-UT; three existing splits each & Evaluate both tasks in their existing cohorts. \\
Patient support & Balanced concentrated and balanced spread & Change patient breadth at matched class patch counts. \\
Readouts & Original CE MLP, logistic regression, cosine \(k\)-NN & Compare learned nonlinear, linear, and local decisions. \\
Evidence & Same frozen features; each allocation's exact patches & Give every readout the same labelled training evidence. \\
Evaluation & Unchanged natural validation and test patients & Measure generalization to unseen patients within each split. \\
\bottomrule
\end{tabularx}
\caption{Patient coverage and classifier choice are varied while each support allocation supplies identical training evidence to all readouts.}
\label{tab:design}
\end{table}

\noindent The historical comparison retains \num{13982}--\num{27202} training patches per BRACS split, with \num{78}--\num{89} concentrated patients versus \num{103} spread patients. TCGA-UT retains \num{36988}--\num{52204} patches, with \num{671}--\num{917} concentrated patients versus \num{4983} spread patients \citep{queisler2026benchmarkResults}. Different shortage doses prevent interpreting a larger effect on one dataset as intrinsic dataset susceptibility.

The spread condition supplies additional patient identities from the original training partition and provides a reference for performance with broader coverage. It changes patient identities and within-patient patch contributions together, so its effect concerns the full coverage intervention, including patient composition and influence.
\section{Readouts}
\label{sec:readouts}
Every readout receives the same training patches within each support condition. Model fitting uses training patients, hyperparameter selection uses validation patients, and final evaluation uses test patients.

\subsection{Original classifier}
The reference is the benchmark's two-layer MLP, trained with cross-entropy loss: a 512-unit hidden layer, ReLU activation, dropout, and a linear class output. The comparison uses the five existing confirmation initializations and their selected configurations, matched to the same features, labels, allocations, and splits. Performance is the mean of the five run-level accuracies, so it describes an individual training procedure averaged over initialization.

\subsection{Linear probe}
Multinomial logistic regression learns a linear decision rule from the original feature coordinates. For training patch \(j\), let \(z_j\in\mathbb{R}^{2560}\) denote its feature vector and \(y_j\) its class. With \(N\) training patches and \(K\) classes, the weights \(W\in\mathbb{R}^{K\times2560}\) and intercept \(b\in\mathbb{R}^{K}\) minimize
\begin{equation}
\mathcal{L}_{\lambda}(W,b)
= -\frac{1}{N}\sum_{j=1}^{N}
\log\left[\operatorname{softmax}(Wz_j+b)\right]_{y_j}
+\frac{\lambda}{2}\lVert W\rVert_F^2.
\label{eq:linear}
\end{equation}
\noindent The first term is the mean cross-entropy loss; the second penalizes large weights, with strength \(\lambda\), using the squared Frobenius norm \(\lVert W\rVert_F^2\), the sum of squared weight entries. The intercept is unpenalized. Features retain their original scale, and training patches receive equal weight. The regularization grid is
\begin{equation}
\lambda\in\{10^{-6},10^{-5},10^{-4},10^{-3},10^{-2},10^{-1},1,10,10^{2}\}.
\label{eq:lambda}
\end{equation}
\noindent A deterministic full-batch solver fits the convex objective to numerical convergence. Prediction selects the class with the largest fitted logit. This readout tests whether a regularized linear rule generalizes better than the original MLP.

\subsection{Nearest neighbours}
Cosine \(k\)-nearest neighbours (\(k\)-NN) predicts a patch's class from the labels of similar training patches. Each feature is normalized to unit length, \(u=z/\lVert z\rVert_2\). Similarity to training patch \(j\) is then the inner product \(u^{\mathsf T}u_j\). If \(\mathcal{N}_k(u)\) denotes the \(k\) most similar training patches, prediction uses an equal vote:
\begin{equation}
\widehat y(u)=\arg\max_{c\in\{1,\ldots,K\}}
\sum_{j\in\mathcal{N}_k(u)}\mathbf{1}\{y_j=c\},
\qquad k\in\{1,5,20,100,250,500\}.
\label{eq:knn}
\end{equation}
\noindent Here \(\widehat y(u)\) is the predicted class, and the indicator \(\mathbf{1}\{y_j=c\}\) equals one when neighbour \(j\) belongs to class \(c\) and zero otherwise. Neighbours are found by exact cosine search. Equal similarities are ordered by patch identifier; class ties for both probes follow a fixed class order. Several neighbours can belong to the same patient, so the vote reflects local patch density.

For each alternative readout, one hyperparameter is selected separately for each dataset and support condition by maximizing validation patient-macro balanced accuracy, averaged equally across the three splits. Numerical ties favour larger \(\lambda\) or \(k\). Test predictions use the selected setting with each split's training-only fit or neighbour bank. The MLP retains its historical hyperparameter search, so gains compare complete fitted classifiers, including their regularization and selection procedures.
\section{Endpoints}
\label{sec:analysis}
Evaluation asks three questions: does another classifier improve on the MLP under shortage, how much does broader patient coverage help, and is the alternative classifier's advantage especially large under shortage? All three use the same accuracy measure, which gives equal importance to patients within each class and then to the classes.

\subsection{Accuracy on unseen patients}
The primary endpoint is \emph{patient-macro balanced accuracy}. Its calculation has three steps. First, for each patient and class, the fraction of correctly classified patches is calculated. Second, these fractions are averaged across patients with that class. Third, the class averages receive equal weight. Thus a patient with many patches cannot dominate a class's score, and a common class cannot dominate the overall score.

For split \(r\), class \(c\), and patient \(i\), let \(\mathcal{J}_{irc}\) be the set of evaluated patches from that patient whose true class is \(c\). The patient's recall for that class is
\begin{equation}
q_{irc}=\frac{1}{|\mathcal{J}_{irc}|}
\sum_{j\in\mathcal{J}_{irc}}\mathbf{1}\{\widehat y_j=c\}.
\label{eq:patient-recall}
\end{equation}
\noindent Here \(|\mathcal{J}_{irc}|\) is the number of those patches, \(\widehat y_j\) is the predicted class of patch \(j\), and the indicator counts correct predictions. Let \(\mathcal{P}_{rc}\) be the patients with at least one evaluated patch of class \(c\). Their recalls are averaged within each class and then across the \(K\) classes:
\begin{equation}
A_r=\frac{1}{K}\sum_{c=1}^{K}
\left(\frac{1}{|\mathcal{P}_{rc}|}\sum_{i\in\mathcal{P}_{rc}}q_{irc}\right).
\label{eq:ba}
\end{equation}
\noindent The inner mean gives each eligible patient equal weight; the outer mean gives each class weight \(1/K\). A patient can contribute to several classes, once within each class. Validation selection applies the same measure to validation patients. Secondary outcomes are the individual class recalls and patch-micro balanced accuracy, which pools patches within each class before averaging class recalls.

\subsection{Classifier gain, coverage benefit, and interaction}
For one dataset, \(A_{h,s}\) denotes patient-macro balanced accuracy for classifier \(h\) under support condition \(s\), averaged equally over the three test splits. The conditions are concentrated (\(\conc\)) and spread (\(\spr\)); the MLP's five initialization-level accuracies are averaged within each split first. The three contrasts are
\begin{align}
G_{h,s}&=A_{h,s}-A_{\mathrm{MLP},s},\label{eq:gain}\\
B_h&=A_{h,\spr}-A_{h,\conc},\label{eq:breadth}\\
I_h&=G_{h,\conc}-G_{h,\spr}=B_{\mathrm{MLP}}-B_h.
\label{eq:interaction}
\end{align}
\noindent Classifier gain \(G_{h,s}\) compares an alternative readout with the MLP while holding support fixed. The central contrast, \(G_{h,\conc}\), asks whether the shortage data support better predictions than the MLP achieves. Coverage benefit \(B_h\) holds the classifier fixed and measures the improvement from broader patient coverage. Interaction \(I_h\) compares the classifier gains across support conditions: a positive value means that the alternative's advantage is larger under shortage. All differences are expressed in percentage points, obtained by multiplying differences in accuracy proportions by 100.

\begin{table}[htbp]
\centering
\small
\renewcommand{\arraystretch}{1.18}
\begin{tabularx}{\textwidth}{@{}p{0.20\textwidth}XX@{}}
\toprule
Contrast & Question & Illustrative calculation \\
\midrule
Classifier gain \(G_{h,\conc}\) & Does another classifier use the shortage data better? & Under concentrated support, MLP accuracy of \(63\%\) and linear accuracy of \(67\%\) give a gain of \(4\) points. \\
Coverage benefit \(B_h\) & How much does this classifier gain from more patients? & Linear accuracy of \(65\%\) under concentrated and \(70\%\) under spread support gives a benefit of \(5\) points. \\
Interaction \(I_h\) & Is the classifier advantage larger under shortage? & A gain over the MLP of \(5\) points under concentrated support and \(1\) point under spread support gives an interaction of \(4\) points. \\
\bottomrule
\end{tabularx}
\caption{The three contrasts separate classifier choice, patient coverage, and their interaction. Numbers are independent hypothetical examples, not experimental results.}
\label{tab:contrasts}
\end{table}

\noindent A positive gain shows usable predictive structure under the same shortage allocation. Positive coverage benefits across all readouts support a role for broader patient coverage. An interaction needs both absolute accuracies and classifier gains for interpretation: a poorly performing probe can have a small coverage benefit, and therefore a positive interaction, without being a useful alternative.

\subsection{Uncertainty across patients and initializations}
Patches from the same patient share biological and acquisition characteristics, so their prediction errors can be correlated. Treating every patch as an independent observation would overstate the amount of independent evidence. Uncertainty is therefore assessed at the patient level.

A \emph{bootstrap} approximates sampling variation by repeatedly resampling or reweighting the observed data and recalculating the statistic of interest. The \emph{Bayesian bootstrap} assigns random positive weights to observations instead of drawing a new sample with replacement. In each of 2,000 replicates here, the weights for the distinct test patients within a dataset are drawn jointly from a \(\mathrm{Dirichlet}(1,\ldots,1)\) distribution. These weights sum to one and treat patients symmetrically before the draw.

Each patient's weight is shared across classifiers, support conditions, classes, and appearances in different test splits, following the benchmark's paired analysis \citep{queisler2026benchmarkProtocol}. The within-class patient mean in Equation~\eqref{eq:ba} becomes a weighted mean, normalized by the total weight of patients present in that class and split. Class and split averages remain equally weighted. Every contrast is then recalculated using the same patient weights on both sides. This pairing preserves the correlation between predictions for the same patients when estimating uncertainty in their differences.

Each replicate also samples five initialization indices with replacement from the MLP's five confirmation runs. The same draw is shared across support conditions, splits, and contrasts involving the MLP; deterministic probe predictions remain fixed. The 2.5th and 97.5th percentiles of the resulting distributions define the reported 95\% intervals. BRACS and TCGA-UT are evaluated separately, with split-level estimates complementing the pooled estimates.
One percentage point is the prespecified practical-relevance threshold for the three signed contrasts. A point estimate above one point is practically promising; a lower interval bound above one point supports an improvement exceeding that threshold. An interval containing zero does not establish equivalence. All intervals below are individual, exploratory 95\% bootstrap intervals.

\section{Results}
\label{sec:results}
Changing the classifier provides little improvement under patient shortage, whereas broader coverage retains a substantial benefit for the learned classifiers. Tables~\ref{tab:results-accuracy} and~\ref{tab:results-contrasts} report the pooled patient-macro estimates and paired contrasts. All primary MLP estimates average the five confirmation runs within each split and then the three splits.

\subsection{Realized comparison and selected settings}
The completed integrity audits passed for both datasets and all six split--support cells per dataset. They verified patient-disjoint partitions, finite nonzero training features, and complete five-run MLP reference blocks. Class-specific patch counts match exactly between concentrated and spread support within every split. Table~\ref{tab:realized-support} gives the realized training allocations. Each BRACS split contains 24 validation and 24 test patients; each TCGA-UT split contains \num{1067} in each partition. These are per-split counts, not counts of distinct patients across the repeated splits.

\begin{table}[htbp]
\centering\small
\begin{tabular}{@{}llrrr@{}}
\toprule
Dataset & Split & Training patches & Patients C & Patients S \\
\midrule
BRACS & 0 & \num{16796} & 85 & 103 \\
 & 1 & \num{27202} & 89 & 103 \\
 & 2 & \num{13982} & 78 & 103 \\
\addlinespace
TCGA-UT & 0 & \num{44792} & 917 & \num{4983} \\
 & 1 & \num{52204} & 752 & \num{4983} \\
 & 2 & \num{36988} & 671 & \num{4983} \\
\bottomrule
\end{tabular}
\caption{Realized training evidence. Patch totals are identical in concentrated (C) and spread (S) support; patient breadth differs. Split identifiers follow the fixed benchmark partitions.}
\label{tab:realized-support}
\end{table}

\noindent Every candidate in both hyperparameter grids was available for selection, and none of the eight selected settings lay at a grid boundary. BRACS selected $\lambda=1$ under concentrated support and $\lambda=0.1$ under spread support, with $k=250$ in both conditions. TCGA-UT selected $\lambda=0.01$ and $k=100$ in both conditions. These settings were selected using validation accuracy only.

\subsection{Classifier gains and remaining coverage benefits}
On BRACS, logistic regression reaches 44.99\% accuracy under concentrated support, close to the MLP's 44.81\%. Its gain of 0.19 points has an interval spanning $-0.59$ to 1.01 points, leaving both its direction and practical relevance uncertain. Under spread support, logistic regression also remains close to the MLP (47.43\% versus 47.72\%). Broader coverage improves the MLP by 2.91 points and logistic regression by 2.44 points; both intervals exclude zero. Neither lower bound exceeds the one-point practical threshold. The linear-probe interaction of 0.47 points is uncertain, so the comparison does not establish a shortage-specific advantage.

\begin{table}[htbp]
\centering\small
\setlength{\tabcolsep}{4pt}
\begin{tabular}{@{}llrrr@{}}
\toprule
Dataset & Readout & Accuracy C (\%) & Accuracy S (\%) & Benefit (points) \\
\midrule
BRACS & MLP & 44.81 & 47.72 & 2.91 \\
 & & [41.67, 48.14] & [43.94, 51.48] & [0.91, 5.22] \\
 & Logistic & 44.99 & 47.43 & 2.44 \\
 & & [41.89, 48.13] & [43.55, 51.18] & [0.31, 4.98] \\
 & $k$-NN & 42.47 & 43.72 & 1.25 \\
 & & [39.26, 45.83] & [39.89, 47.29] & [$-0.65$, 3.42] \\
\addlinespace
TCGA-UT & MLP & 74.66 & 79.43 & 4.76 \\
 & & [73.28, 75.91] & [78.15, 80.61] & [4.12, 5.40] \\
 & Logistic & 75.15 & 79.84 & 4.70 \\
 & & [73.77, 76.39] & [78.50, 81.08] & [4.04, 5.34] \\
 & $k$-NN & 68.09 & 71.17 & 3.08 \\
 & & [66.85, 69.35] & [69.89, 72.41] & [2.56, 3.58] \\
\bottomrule
\end{tabular}
\caption{Patient-macro balanced accuracy and paired coverage benefit $B_h=A_{h,\spr}-A_{h,\conc}$. Brackets give 95\% intervals beneath point estimates. Coverage benefits remain positive for the MLP and logistic regression on both datasets; the BRACS nearest-neighbour interval includes zero.}
\label{tab:results-accuracy}
\end{table}

\noindent On TCGA-UT, logistic regression improves on the MLP by 0.48 points under concentrated support (95\% interval: 0.28 to 0.70) and 0.42 points under spread support (0.17 to 0.67). Both gains exclude zero, but their upper bounds remain below one point. Broader coverage still adds 4.70 points to logistic regression, close to the MLP's 4.76-point benefit. The interaction is only 0.07 points ($-0.23$ to 0.36), which does not support a practically important shortage-specific gain. The concentrated linear probe also remains 4.28 points below the spread MLP by point estimate; no paired interval for this cross-condition comparison is included in the available results.

\begin{table}[htbp]
\centering\small
\setlength{\tabcolsep}{4pt}
\begin{tabular}{@{}llrrr@{}}
\toprule
Dataset & Probe & Shortage gain & Spread gain & Interaction \\
\midrule
BRACS & Logistic & 0.19 & $-0.28$ & 0.47 \\
 & & [$-0.59$, 1.01] & [$-1.10$, 0.53] & [$-0.68$, 1.69] \\
 & $k$-NN & $-2.34$ & $-4.00$ & 1.66 \\
 & & [$-3.99$, $-0.56$] & [$-6.09$, $-1.76$] & [0.24, 3.04] \\
\addlinespace
TCGA-UT & Logistic & 0.48 & 0.42 & 0.07 \\
 & & [0.28, 0.70] & [0.17, 0.67] & [$-0.23$, 0.36] \\
 & $k$-NN & $-6.57$ & $-8.25$ & 1.68 \\
 & & [$-7.25$, $-5.84$] & [$-8.99$, $-7.51$] & [1.03, 2.33] \\
\bottomrule
\end{tabular}
\caption{Paired classifier gains over the MLP and their interaction, in percentage points, with 95\% intervals. Positive interactions for $k$-NN occur alongside negative gains: they indicate a smaller coverage benefit, not successful shortage mitigation.}
\label{tab:results-contrasts}
\end{table}

\noindent Nearest neighbours perform worse than the MLP under both conditions on both datasets, with all four gain intervals wholly below zero. Their shortage deficits are 2.34 points on BRACS and 6.57 points on TCGA-UT. Spread support increases nearest-neighbour accuracy by 1.25 points on BRACS, with an interval including zero, and by 3.08 points on TCGA-UT, with an interval wholly above one point. The resulting positive interactions therefore reflect reduced responsiveness to coverage rather than an advantage over the reference classifier.

\subsection{Split variation and secondary outcomes}
The deterministic probes permit a direct check of split variation (Table~\ref{tab:probe-splits}). Logistic regression benefits from spread support in all three splits of each dataset. Its BRACS benefit varies from 0.27 to 4.60 points, compared with 4.31 to 5.10 points on TCGA-UT. BRACS nearest neighbours lose about 0.75 points in splits 0 and 1 and gain 5.25 points in split 2. Their pooled positive estimate thus conceals inconsistent responses. TCGA-UT nearest neighbours benefit in every split, by 2.33 to 4.31 points.

\begin{table}[htbp]
\centering\small
\begin{tabular}{@{}llrrrr@{}}
\toprule
 & & \multicolumn{2}{c}{Logistic regression} & \multicolumn{2}{c}{$k$-NN} \\
\cmidrule(lr){3-4}\cmidrule(l){5-6}
Dataset & Split & C & S & C & S \\
\midrule
BRACS & 0 & 47.37 & 49.82 & 45.41 & 44.67 \\
 & 1 & 45.80 & 46.08 & 43.87 & 43.11 \\
 & 2 & 41.80 & 46.40 & 38.13 & 43.38 \\
\addlinespace
TCGA-UT & 0 & 75.58 & 79.90 & 68.79 & 71.12 \\
 & 1 & 74.50 & 79.18 & 67.25 & 71.56 \\
 & 2 & 75.36 & 80.46 & 68.22 & 70.84 \\
\bottomrule
\end{tabular}
\caption{Split-level patient-macro balanced accuracy (\%) for the selected deterministic probes. These descriptive estimates are correlated across splits and do not constitute independent replications.}
\label{tab:probe-splits}
\end{table}

\noindent Patch-micro balanced accuracy changes the BRACS pattern. Averaged over splits, logistic regression scores 50.21\% under concentrated and 50.04\% under spread support, a difference of $-0.17$ points. Nearest neighbours score 44.72\% and 42.98\%, a difference of $-1.74$ points. Thus the positive patient-macro coverage estimates on BRACS do not extend to an endpoint that gives more influence to patients contributing more patches within a class. On TCGA-UT, the endpoints agree in direction: patch-micro accuracy increases from 75.77\% to 80.07\% for logistic regression and from 68.67\% to 71.86\% for nearest neighbours. These secondary estimates have no reported uncertainty intervals.

Across the equally averaged split-level class recalls, spread support improves 5 of 7 BRACS classes for each probe, 25 of 30 TCGA-UT classes for logistic regression, and 26 of 30 for nearest neighbours. These descriptive counts show that coverage benefits are not confined to a single class, while also showing that not every class benefits. They neither measure the size of each class's effect nor establish class-specific improvements with uncertainty control.

The available split-level and secondary MLP summaries represent only the first confirmation run, whereas the primary pooled analysis uses all five. They are therefore excluded from these secondary comparisons to preserve the stated MLP estimand. This restriction does not affect the primary accuracies, classifier gains, coverage benefits, or interactions in Tables~\ref{tab:results-accuracy} and~\ref{tab:results-contrasts}.


\section{Discussion}
\label{sec:interpretation}
The principal finding is that replacing the MLP does not remove the performance cost of concentrated patient support. Logistic regression comes closest to the MLP in both datasets, but its shortage gain is small relative to its own benefit from broader coverage. On TCGA-UT, the linear classifier improves over the MLP under both support conditions, demonstrating a modest classifier limitation alongside a substantially larger coverage effect. The near-zero interaction does not identify this gain as specific to shortage. On BRACS, the linear classifier supplies no clear gain over the MLP, whereas broader coverage improves both learned classifiers.

The evidence for coverage is strongest on TCGA-UT: all three readouts benefit, with lower interval bounds above one percentage point. BRACS supports a narrower conclusion. Its MLP and linear-probe coverage intervals exclude zero but do not establish benefits greater than one point; the nearest-neighbour interval includes zero and its split-level effects conflict. Thus the results support broader patient coverage across both tasks without claiming equally consistent evidence across all readouts and datasets.

Nearest neighbours illustrate why the interaction cannot be interpreted alone. Their positive interaction reflects a smaller spread benefit than the MLP's, while their absolute accuracy is lower under both support conditions. Even the TCGA-UT interaction, whose interval lies above the practical threshold, therefore does not identify a useful shortage remedy. A classifier must first offer a favourable shortage gain before a positive interaction supports its use for this purpose.

The coverage intervention changes patient identities and their patch contributions together. Its benefit therefore supports the value of the broader labelled allocation, rather than isolating a causal effect of patient count alone. Likewise, logistic regression and the MLP differ in fitting and regularization, so the small TCGA-UT gain cannot be attributed solely to linear versus nonlinear decision geometry. Stronger regularization selected for concentrated BRACS support is consistent with a need to constrain fitting with fewer patients, but this selection pattern does not establish the mechanism of the coverage deficit.

These findings motivate a follow-up experiment asking whether a small but deliberately diverse patient set can retain the benefit of broader coverage. At several training-patient budgets, random patient selection could be compared with selection that spreads the chosen patients across the frozen feature space, using training data only. Class-specific patch totals would remain fixed, with matched patient counts and per-patient contribution limits between selection strategies. Repeated patient draws, evaluated with the MLP and logistic regression on the fixed held-out patients, would measure both accuracy and sensitivity to patient composition. An advantage for diversity-based selection would suggest that which patients are labelled matters beyond their number; a persistent benefit from larger patient sets under both strategies would strengthen the case for acquiring broader labelled coverage.

\section{Limitations}
\label{sec:budget}
The comparison concerns fixed allocations and cohorts whose benchmark outcomes motivated the question. The intervals describe held-out patient and MLP initialization variation conditional on these allocations, selected hyperparameters, and existing splits. They omit variation from new training-patient samples and hyperparameter selection. Repeated splits share patients, and paired bootstrap weights only account for part of the dependence: overlapping training and validation data remain a limitation. Fresh patient draws or an independent held-out cohort would strengthen confirmation.

The MLP and alternative readouts differ in their fitting and hyperparameter searches as well as their decision rules. A gain therefore identifies a better-performing classification procedure; attributing it specifically to architecture, optimization, or regularization requires further comparisons. All contrasts are exploratory, and the intervals describe individual effects without simultaneous coverage across comparisons.

The findings concern patch classification with one frozen encoder and the tested patient-coverage interventions. A useful shortage-trained probe establishes available predictive structure, while persistent spread benefits establish value from broader labelled coverage for the tested readouts. Failure to find a gain leaves open the possibility that other classifiers or representations could perform better.
\bibliographystyle{plainnat}
\bibliography{references}
\end{document}
